\chapter{Contrats à terme de taux à court terme}\label{ch:m2:short-term-interest-rate-futures}

Un trader regarde le contrat à terme SOFR à un mois de décembre 2026. Il se négocie à 96.040,
ce qui signifie que le taux au jour le jour sera en moyenne de 3.960~\% sur le mois. La
Federal Reserve se réunit le 9 décembre~; le contrat de novembre indique que le taux
sera de 3.8975~\% jusqu'à la réunion. Neuf jours de décembre à 3.8975~\% et
vingt-deux jours à un taux inconnu doivent donner une moyenne de 3.960~\%~: le taux inconnu est
3.986~\%, soit 8.8 points de base de plus, et si les seuls choix sont de maintenir ou de
relever d'un quart de point, le marché intègre une probabilité de 35~\% d'une hausse. Le
desk du trader pense que la probabilité est plus faible, mais avant de vendre le contrat il
doit se poser une question de plus, sur le dernier jour de décembre, où le
taux au jour le jour est généralement plus élevé pour des raisons qui n'ont rien à voir avec la
Federal Reserve. Les prix de ce chapitre sont illustratifs~; les dates de réunion
et les règles des contrats sont réelles. Il montre comment les \omterm{def:m2:short-term-interest-rate-futures:future}{contrats à terme sur taux au jour le jour}
transforment le calendrier d'une banque centrale en prix, et comment les relire.

\section{Contrats à terme sur taux au jour le jour}

\begin{definition}[Contrat à terme sur taux au jour le jour]\label{def:m2:short-term-interest-rate-futures:future}
Un \emph{contrat à terme sur taux au jour le jour}\index{contrat à terme sur taux au jour le jour} est un contrat à terme réglé en espèces
dont le prix final vaut 100 moins un taux calculé à partir des fixings
quotidiens d'un indice de référence au jour le jour (\cref{def:m2:central-banks-and-the-short-rate:rfr})
sur une période de référence~: leur moyenne arithmétique sur un mois civil pour
les contrats à un mois, leur taux composé sur un trimestre pour les contrats à trois mois. Avant l'échéance, le prix est l'espérance de marché de ce
contrats. Avant l'échéance, le prix est l'espérance de marché de ce
règlement-livraison, corrigé de la convexité (\cref{sec:m2:short-term-interest-rate-futures:convexity}).
\end{definition}

Coter $100 - R$ fait que le contrat se comporte comme une obligation~: il monte quand les taux
baissent. Chaque contrat a une valeur fixe par point de base, si bien que le
risque de taux d'une position est simplement égal à son nombre de contrats multiplié par cette valeur, et
les gains et les pertes sont versés chaque jour sous forme de marge de variation (livre~1,
chapitre~5).

\begin{dated}{2026-09}{Contrats sur taux au jour le jour}\label{dat:m2:short-term-interest-rate-futures:contracts}
\textbf{CME three-month SOFR} (\texttt{SR3})~: 2\,500~USD fois l'indice, 25~USD par
point de base~; 100 moins le SOFR composé sur le \omterm{def:m2:short-term-interest-rate-futures:imm}{trimestre de référence}, qui s'étend
du troisième mercredi du troisième mois précédant le mois de livraison jusqu'au
troisième mercredi du mois de livraison~; pas de cotation d'un quart de point de base pour les
contrats les plus proches, la moitié sinon~; 39 contrats trimestriels, soit près de dix ans.
\textbf{CME one-month SOFR} (\texttt{SR1})~: 4\,167~USD fois l'indice, 41.67~USD par
point de base~; 100 moins la moyenne arithmétique du SOFR sur le mois civil.
\textbf{ICE three-month SONIA}~: 2\,500~GBP fois l'indice~; SONIA composé à partir du
troisième mercredi du mois de livraison jusqu'au trimestre suivant~; 25
contrats trimestriels. \textbf{Euribor}, le taux interbancaire de la zone euro encore en vigueur,
est publié à 11:00 CET pour cinq maturités par son administrateur, et les contrats à terme
sur l'Euribor à trois mois se négocient aux côtés de ceux sur taux au jour le jour.
\end{dated}

\section{Dates IMM, packs, bundles et strips}

\begin{definition}[Date IMM et trimestre de référence]\label{def:m2:short-term-interest-rate-futures:imm}
Une \emph{date IMM}\index{date IMM} est le troisième mercredi de mars, juin,
septembre ou décembre. Le \emph{trimestre
de référence}\index{trimestre de référence} d'un contrat à trois mois est la période entre deux
dates IMM consécutives sur laquelle son taux est composé.
\end{definition}

Ces dates servent au-delà des contrats à terme~: les accords de taux à terme (FRA) et les swaps sont
souvent fixés à ces dates, de sorte que leurs couvertures s'alignent sur le \omterm{def:m2:short-term-interest-rate-futures:strip}{strip de contrats à terme}.

\begin{definition}[Strip de contrats à terme, pack, bundle]\label{def:m2:short-term-interest-rate-futures:strip}
Un \emph{strip de contrats à terme}\index{strip de contrats à terme} est une suite de
contrats trimestriels consécutifs, dont les taux décrivent ensemble la trajectoire du taux court.
Un \emph{pack}\index{pack} est un ensemble de quatre contrats trimestriels consécutifs négociés en
quantités égales comme un seul instrument, soit une année du strip~; un
\emph{bundle}\index{bundle} est l'ensemble des $n$ premières années de contrats consécutifs
négociés ensemble, depuis l'échéance la plus proche.
\end{definition}

\begin{proposition}[Un strip est un taux à terme]\label{prop:m2:short-term-interest-rate-futures:strip}
Si des périodes de référence consécutives de longueurs $\delta_1, \dots, \delta_n$ (en
années) ont pour taux de contrats à terme $R_1, \dots, R_n$ (convexité négligée), un dépôt
couvrant toutes ces périodes a pour taux
\[
R_{1\dots n} \;=\; \frac{\prod_{i}(1 + R_i\,\delta_i) - 1}{\sum_i \delta_i},
\]
et l'achat du strip fixe ce taux pour un investisseur qui reconduira un
dépôt sur les périodes.
\end{proposition}
\begin{proof}
Une unité placée au taux de chaque période croît jusqu'au produit~; les
règlements-livraisons quotidiens du strip reproduisent l'écart entre les taux réalisés et les
taux verrouillés, période par période, au premier ordre.
\end{proof}

\begin{omfigure}
\begin{tikzpicture}[font=\footnotesize,x=1.05cm]
\draw[thick,-{Stealth[length=2.5mm]}] (0,0) -- (12.9,0) node[right]{temps};
\foreach \x/\l in {0/1 sept.,1/1 oct.,2/1 nov.,3/1 déc.,4/1 janv.,5/1 févr.,6/1 mars}{\draw (\x*1.9,0.1) -- (\x*1.9,-0.1) node[below]{\l};}
% one-month contracts
\foreach \x/\l in {1/oct.,2/nov.,3/déc.,4/janv.}{\draw[omDef,line width=2.5pt] (\x*1.9+0.05,0.45) -- (\x*1.9+1.85,0.45); \node[above,text=omDef] at (\x*1.9+0.95,0.5) {SR1 \l};}
% x = 1.9 x (months after 1 Sep + day fraction): 16 Sep 0.50, 16 Dec 3.48, 17 Mar 6.52
\draw[omThm,line width=2.5pt] (0.50*1.9,1.35) -- (3.48*1.9,1.35);
\node[above,text=omThm] at (2.0*1.9,1.4) {SR3 décembre~: 16 sept. -- 16 déc.};
\draw[omProp,line width=2.5pt] (3.48*1.9,1.35) -- (6.52*1.9,1.35);
\node[above,text=omProp] at (5.0*1.9,1.4) {SR3 mars~: 16 déc. -- 17 mars};
% FOMC decisions
\foreach \x/\l in {0.50/16 sept.,1.87/28 oct.,3.26/9 déc.,4.84/27 janv.}{\draw[omThm,thick,dashed] (\x*1.9,-0.6) -- (\x*1.9,0.3); \node[below,text=omThm] at (\x*1.9,-0.6) {\l};}
\node[anchor=east,text=omThm] at (-0.1,-0.8) {FOMC:};
\end{tikzpicture}

\omcaption{Périodes de référence des contrats de ce chapitre, septembre 2026
à mars 2027. Les contrats à un mois font la moyenne de leur mois civil~; les contrats à trois mois
composent leur taux de \omterm{def:m2:short-term-interest-rate-futures:imm}{date IMM} à \omterm{def:m2:short-term-interest-rate-futures:imm}{date IMM}, et sont nommés d'après le mois
où leur période \emph{se termine}. Les lignes en pointillés sont les dates de décision du Federal Open Market
Committee~: ce sont les jours de décision~; chaque nouveau taux s'applique à partir du
lendemain.}\label{fig:m2:short-term-interest-rate-futures:calendar}
\end{omfigure}

\section{Lire une réunion}

\begin{definition}[Trajectoire implicite des taux directeurs]\label{def:m2:short-term-interest-rate-futures:path}
La \emph{trajectoire implicite des taux directeurs}\index{trajectoire implicite des taux directeurs} est la suite des taux directeurs que le marché anticipe après chaque réunion programmée de la banque centrale, lue
à partir des cours de contrats à terme sur taux à court terme
(ou des swaps de taux au jour le jour, OIS), comme si ces cours étaient des anticipations,
la prime de risque étant négligée.
\end{definition}

\begin{method}[Lire une réunion à partir d'un contrat à un mois]\label{met:m2:short-term-interest-rate-futures:meeting}
Soit un mois de $N$ jours civils dont la décision tombe le jour $m$, le nouveau taux
s'appliquant à partir du jour $m + 1$. Soit $\bar R$ le taux implicite
dans le contrat du mois, $100 - F$, et $r_0$ le taux anticipé avant la réunion (à partir
du contrat du mois précédent ou du fixing courant). On retire toute prime de \omterm{def:m2:money-markets:turn}{passage de fin de période}
connue $\tau$ sur $k$ jours, puis
\[
r_1 \;=\; \frac{N\bar R - k\tau - m\,r_0}{N - m},
\qquad
p \;=\; \frac{r_1 - r_0}{\Delta},
\]
où $p$ est la probabilité implicite d'un mouvement de taille $\Delta$ si les seules
issues sont l'absence de mouvement et $\Delta$. Il faut enchaîner les mois pour obtenir le taux
anticipé après chaque réunion.
\end{method}

\begin{example}[Trois réunions]\label{ex:m2:short-term-interest-rate-futures:path}
À partir d'un taux de 3.87~\% après la décision de septembre 2026, avec des cours illustratifs
de 96.1275 (octobre), 96.1025 (novembre), 96.040 (décembre) et 95.990
(janvier 2027)~: le contrat de novembre implique 3.8975~\% après la réunion d'octobre
soit une probabilité de 11~\% d'une hausse~; décembre implique 3.9856~\% après la réunion du 9
décembre, soit 35~\% de probabilité~; janvier implique 4.175~\% après la réunion du 27 janvier
(\cref{fig:m2:short-term-interest-rate-futures:path}).
\end{example}

\begin{remark}[Les réunions tardives sont fragiles]\label{rem:m2:short-term-interest-rate-futures:fragile}
Le multiplicateur $N/(N - m)$ transforme une erreur de cours en erreur sur $r_1$. La
réunion de janvier tombe le 27, si bien que seuls quatre jours du mois portent
le nouveau taux~: une erreur d'un point de base sur le cours de janvier déplace le taux
implicite après la réunion de $31/4 = 7.75$ points de base, soit un tiers d'une hausse. La réunion d'octobre
est pire (trois jours). Pour les réunions tardives, les desks utilisent le contrat du mois suivant,
dans lequel le nouveau taux s'applique tout le mois, comme l'exemple l'a fait pour
octobre.
\end{remark}

\begin{omfigure}
\begin{tikzpicture}
\begin{axis}[width=11cm,height=4.8cm,
  ylabel={taux au jour le jour implicite (\%)},
  tick label style={font=\small},label style={font=\small},
  xmin=10,xmax=153,ymin=3.82,ymax=4.22,grid=major,grid style={gray!25},
  xtick={30,61,91,122,153},xticklabels={1 oct.,1 nov.,1 déc.,1 janv.,1 févr.},
  table/col sep=comma]
\addplot[omDef,very thick,const plot] table[x=t,y=rate]{figdata/markets-2/08-short-term-interest-rate-futures/path.csv};
\draw[omThm,thick,dashed] (axis cs:57,3.82) -- (axis cs:57,4.22);
\draw[omThm,thick,dashed] (axis cs:99,3.82) -- (axis cs:99,4.22);
\draw[omThm,thick,dashed] (axis cs:148,3.82) -- (axis cs:148,4.22);
\node[font=\footnotesize,anchor=north west,text=omThm] at (axis cs:101,3.975) {9 déc.~: 35~\%};
\node[font=\footnotesize,anchor=east,text=omThm] at (axis cs:146,4.12) {27 janv.};
\end{axis}
\end{tikzpicture}

\omcaption{La trajectoire du taux au jour le jour implicite dans les contrats à un mois illustratifs
de l'\cref{ex:m2:short-term-interest-rate-futures:path}, avec un saut à
chaque décision (en pointillés). Ces sauts sont des anticipations, moyennes entre maintien et
hausse, et non des prévisions d'un mouvement de un quart de point~: 35~\% de 25 points de base font les 8.8
de décembre. Données~: le tutoriel du chapitre.}\label{fig:m2:short-term-interest-rate-futures:path}
\end{omfigure}

\begin{omfigure}
\begin{tikzpicture}
\begin{axis}[width=11cm,height=4.4cm,
  xlabel={prime de passage de fin d'année supposée au 31 décembre (pb)},ylabel={probabilité de hausse (\%)},
  tick label style={font=\small},label style={font=\small},
  xmin=0,xmax=40,ymin=26,ymax=37,grid=major,grid style={gray!25},
  table/col sep=comma]
\addplot[omThm,very thick] table[x=turn,y=prob]{figdata/markets-2/08-short-term-interest-rate-futures/turn.csv};
\end{axis}
\end{tikzpicture}

\omcaption{La probabilité de hausse de décembre implicite par le même cours croît avec le
passage de fin d'année supposé au fixing du 31 décembre~: 35.2~\% sans passage, 33.4~\%
avec 10 points de base, soit 0.18 point de pourcentage de moins par point de base de passage.
Le \omterm{def:m2:money-markets:turn}{passage de fin de période} est un fait du marché monétaire (voir au \cref{ch:m2:money-markets}) qu'une lecture
de la politique monétaire doit d'abord retirer. Données~: le problème du
week-end du chapitre.}\label{fig:m2:short-term-interest-rate-futures:turn}
\end{omfigure}

\section{L'ajustement de convexité}\label{sec:m2:short-term-interest-rate-futures:convexity}

Un contrat à terme est réglé chaque jour~; un forward rate agreement ou un swap est réglé
une seule fois. La différence compte, car les gains d'une position vendeuse sur contrats à terme
(qui profite de la hausse des taux) arrivent lorsque les taux sont élevés et peuvent être réinvestis
à des taux élevés, et ses pertes arrivent lorsque les taux sont bas et peuvent être financées
à bon marché. Le vendeur en profite, si bien qu'il accepte un taux de contrat à terme
supérieur au taux forward.

\begin{definition}[Ajustement de convexité]\label{def:m2:short-term-interest-rate-futures:convexity}
L'\emph{ajustement de convexité}\index{ajustement de convexité} d'un contrat à terme sur taux court
est l'écart entre son taux implicite et le taux forward pour
la même période.
\end{definition}

\begin{proposition}[L'ajustement dans un modèle gaussien]\label{prop:m2:short-term-interest-rate-futures:holee}
Si le taux court suit $dr = \theta(t)\,dt + \sigma\,dW$ (le modèle de Ho--Lee),
le taux du contrat à terme pour une période $[T_1, T_2]$ dépasse le taux forward d'environ
$\tfrac12\,\sigma^2\,T_1\,T_2$. Le livre~6 le démontre.
\end{proposition}
\admitted

\begin{example}[Cinq et dix ans]\label{ex:m2:short-term-interest-rate-futures:convexity}
Avec une volatilité normale de 1~\% par an, le contrat à trois mois débutant dans cinq ans
porte $\frac12 \times 0.01^2 \times 5 \times 5.25 = 13.1$ points de base
d'ajustement~; celui débutant dans dix ans, 51.25. À l'avant de la courbe, il est négligeable~;
à l'arrière du \omterm{def:m2:short-term-interest-rate-futures:strip}{strip de contrats à terme}, il est plus grand qu'un mouvement habituel dans une journée, et une
courbe construite à partir des contrats à terme sans lui est fausse
(\cref{fig:m2:short-term-interest-rate-futures:convexity}).
\end{example}

\begin{omfigure}
\begin{tikzpicture}
\begin{axis}[width=11cm,height=4.6cm,
  xlabel={début de la période de trois mois (années)},ylabel={ajustement (pb)},
  tick label style={font=\small},label style={font=\small},
  xmin=0,xmax=10,ymin=0,ymax=78,grid=major,grid style={gray!25},
  legend columns=3,legend style={font=\footnotesize,at={(0.5,-0.3)},anchor=north,draw=none,fill=none,/tikz/every even column/.append style={column sep=8pt}},
  table/col sep=comma]
\addplot[omDef,thick,dashed] table[x=t1,y=s08]{figdata/markets-2/08-short-term-interest-rate-futures/convexity.csv};
\addplot[omThm,very thick] table[x=t1,y=s10]{figdata/markets-2/08-short-term-interest-rate-futures/convexity.csv};
\addplot[omProp,thick,dotted] table[x=t1,y=s12]{figdata/markets-2/08-short-term-interest-rate-futures/convexity.csv};
\legend{$\sigma = 0.8\%$,$\sigma = 1.0\%$,$\sigma = 1.2\%$}
\end{axis}
\end{tikzpicture}

\omcaption{L'\omterm{def:m2:short-term-interest-rate-futures:convexity}{ajustement de convexité} de Ho--Lee d'un taux de contrat à terme à trois mois selon
le début de sa période, pour trois volatilités. Il croît avec le carré de l'horizon et
de la volatilité~: négligeable pour les premiers contrats, de dizaines de points de base pour
les derniers. Données~: le tutoriel du chapitre.}\label{fig:m2:short-term-interest-rate-futures:convexity}
\end{omfigure}

\section{Tutoriel~: extraction d'une trajectoire de politique monétaire}\label{tut:m2:short-term-interest-rate-futures:path}

\begin{tutorial}
\textbf{Objectif.} Lire le taux attendu après chaque réunion à partir des contrats à terme à un mois,
avec et sans passage de fin d'année, et calculer les taux de strip et les
\omterm{def:m2:short-term-interest-rate-futures:convexity}{ajustements de convexité}. \textbf{État final~:}
\cref{fig:m2:short-term-interest-rate-futures:path,fig:m2:short-term-interest-rate-futures:turn}
et les valeurs numériques de l'\cref{ex:m2:short-term-interest-rate-futures:path}.
\begin{tutsteps}
\item \textbf{Un mois.} Retirer le \omterm{def:m2:money-markets:turn}{passage de fin de période}, puis résoudre pour le taux après la réunion.
\omcode{code/firm/meetings/firm_meetings.py}{32}{39}{Le taux après une réunion, calculé à partir de la moyenne du mois.}\label{lst:m2:short-term-interest-rate-futures:after}
\item \textbf{Enchaîner les mois}, et convertir les changements de taux en probabilités.
\omcode{code/firm/meetings/firm_meetings.py}{42}{53}{La trajectoire de politique monétaire, et la probabilité d'un mouvement.}\label{lst:m2:short-term-interest-rate-futures:path}
\item \textbf{Strips.}
\omcode{code/firm/meetings/firm_meetings.py}{56}{62}{Le taux d'un strip de périodes consécutives.}\label{lst:m2:short-term-interest-rate-futures:strip}
\item \textbf{Exécuter} \texttt{stir\_demo.path()} et \texttt{stir\_demo.december(10)}~:
3.9856~\% et 35.2~\% sans \omterm{def:m2:money-markets:turn}{passage de fin de période}, 3.9810~\% et 33.4~\% avec 10 points de base.
\end{tutsteps}
\textbf{Pour aller plus loin.} Remplacer les contrats à un mois par des contrats à trois mois
qui se composent sur les trimestres IMM, avec deux réunions dans un même trimestre,
et résoudre pour les deux (l'\cref{exo:m2:short-term-interest-rate-futures:7}).
\end{tutorial}

\section{Construction~: l'extracteur de trajectoire de politique monétaire}\label{bld:m2:short-term-interest-rate-futures:meetings}

\begin{build}
\bfield{Rôle} Les traders de taux de la société miniature, et son constructeur de courbe
(du \cref{ch:m2:interest-rate-swaps}) ont besoin de la trajectoire de politique monétaire attendue par le marché~:
les traders pour prendre position réunion par réunion, la courbe pour placer ses paliers
aux dates de réunion plutôt que de les lisser. \bfield{Interface} \texttt{MonthContract(year, month, price, meeting,
turn\_days, turn\_bp)}~; \texttt{rate\_after(contract, rate\_before)}~;
\texttt{policy\_path(contracts, start\_rate)}~;
\texttt{move\_probability(after, before, step)}~;
\texttt{strip\_rate(rates, fractions)}~;
\texttt{convexity\_adjustment(sigma, t1, t2)}.
\bfield{Règles} Moyennes sur les jours calendaires~; le nouveau taux s'applique à partir du lendemain
d'une décision~; une prime de \omterm{def:m2:money-markets:turn}{passage de fin de période} retirée avant la résolution~; taux en pourcentage.
\bfield{Tests d'acceptation} \texttt{code/firm/meetings/tests/}~: un mois sans réunion renvoie sa moyenne~; une
trajectoire construite est retrouvée exactement~; un \omterm{def:m2:money-markets:turn}{passage de fin de période} est retiré~;
l'enchaînement des mois~; la capitalisation du strip~; la valeur de l'ajustement
à cinq ans.
\bfield{Pour aller plus loin} Contrats à trois mois (capitalisation, plusieurs réunions par trimestre)~; la probabilité
d'un mouvement de 25 contre 50 points de base à partir de deux contrats adjacents~; la série
des trajectoires implicites au cours du temps, comme produit de données.
\end{build}

\begin{omsources}
\item CME Group, Rulebook Chapters 460 and 461 (contrats à terme SOFR à trois mois
et à un mois)~; ICE, \emph{Three Month SONIA Index Futures}~; EMMI, \emph{Euribor}.
\item Federal Reserve, \emph{FOMC meeting calendars}, 2026 et 2027.
\item T.\,S.\,Y.\ Ho et S.-B.\ Lee, ``Term structure movements and pricing
interest rate contingent claims'', \emph{Journal of Finance} 41, 1986.
\item G.\ Burghardt, \emph{The Eurodollar Futures and Options Handbook},
McGraw-Hill, 2003, pour la pratique de marché dont les contrats SOFR ont hérité.
\end{omsources}

\section{Exercices}

\begin{exercise}[$\star$]\label{exo:m2:short-term-interest-rate-futures:1}
Un contrat SOFR à trois mois se négocie à 95.965. Quel taux implique-t-il~? Il monte de
3.5 points de base~: combien de ticks d'un demi-point de base cela représente-t-il, et quel est le gain sur 200
contrats~?
\end{exercise}

\begin{exercise}[$\star$]\label{exo:m2:short-term-interest-rate-futures:2}
Donner les quatre \omterm{def:m2:short-term-interest-rate-futures:imm}{dates IMM} de 2027 et le \omterm{def:m2:short-term-interest-rate-futures:imm}{trimestre de référence} du contrat
SOFR à trois mois de juin 2027.
\end{exercise}

\begin{exercise}[$\star$]\label{exo:m2:short-term-interest-rate-futures:3}
Le contrat à un mois de novembre se négocie à 96.1025. Quel taux moyen implique-t-il~?
Que coûte une hausse de 2 points de base de ce taux à un détenteur de 100
contrats achetés~?
\end{exercise}

\begin{exercise}[$\star\star$]\label{exo:m2:short-term-interest-rate-futures:4}
La réunion d'octobre a lieu le 28. De combien de points de base une erreur de
1 point de base sur le prix d'octobre déplace-t-elle le taux implicite après la réunion~?
Quel contrat utiliser à la place, et pourquoi~?
\end{exercise}

\begin{exercise}[$\star\star$]\label{exo:m2:short-term-interest-rate-futures:5}
Quatre contrats trimestriels consécutifs impliquent 3.95~\%, 4.05~\%, 4.10~\% et 4.12~\% sur
périodes de 91 jours chacune (exact/360). Donner le taux de strip sur un an.
\end{exercise}

\begin{exercise}[$\star\star$]\label{exo:m2:short-term-interest-rate-futures:6}
Un trader couvre une position en livres sterling avec le contrat à terme SONIA à trois mois
«~décembre~», en croyant qu'il couvre la même période que le contrat à terme SOFR à trois
mois «~décembre~». Qu'est-ce qui cloche~?
\end{exercise}

\begin{exercise}[$\star\star\star$]\label{exo:m2:short-term-interest-rate-futures:7}
\emph{Programmation.} Avec \texttt{convexity\_adjustment}, donner l'\omterm{def:m2:short-term-interest-rate-futures:convexity}{ajustement de convexité} du
contrat à trois mois débutant dans cinq ans et dans dix ans, pour une volatilité de 1~\%.
De combien varie-t-il si la volatilité passe à 1.2~\%~?
\end{exercise}

\begin{exercise}[$\star\star\star$]\label{exo:m2:short-term-interest-rate-futures:8}
\emph{Repérer l'erreur.} «~Le contrat d'octobre implique une probabilité de 10~\% d'une hausse
le 28 octobre, et le contrat de novembre 11~\%~: le marché est incohérent en interne,
et la différence constitue un arbitrage.~» Corriger cette affirmation.
\end{exercise}

\section{Problème~: hausse ou maintien}

\begin{problem}[{Problème du week-end --- la réunion de décembre 2026}]\label{pb:m2:short-term-interest-rate-futures:1}
Le contrat à terme SOFR à un mois de novembre se négocie à 96.1025 et celui de décembre à
96.0400. Le Comité décide le 9 décembre~; le nouveau taux s'applique à partir du 10
décembre. Les seules issues possibles sont le maintien et une hausse d'un quart de point.
Chaque contrat vaut 41.67~USD par point de base.

\medskip\noindent\textbf{Partie I --- La lecture.}
\begin{enumerate}
\item Quel taux le marché anticipe-t-il avant la réunion~?
\item Combien de jours de décembre sont au taux ancien, et combien au taux nouveau~?
\item Donner le taux après réunion implicite dans le prix de décembre.
\item Donner la probabilité implicite d'une hausse.
\item Combien de points de base de prix de décembre représente un point de
pourcentage de probabilité~?
\end{enumerate}

\medskip\noindent\textbf{Partie II --- Le \omterm{def:m2:money-markets:turn}{passage de fin de période}.}
\begin{enumerate}[resume]
\item Le fixing du 31 décembre devrait dépasser de 10 points de base celui des autres jours.
Recalculer le taux après réunion et la probabilité.
\item Et pour 20 points de base~?
\item De combien chaque point de base de passage supposé réduit-il la probabilité~?
\item Pourquoi le \omterm{def:m2:money-markets:turn}{passage de fin de période} se situe-t-il dans le contrat de décembre et non dans celui de janvier~?
\item Comment estimer le passage à partir des prix de marché plutôt que de le supposer~?
\end{enumerate}

\medskip\noindent\textbf{Partie III --- La transaction.}
\begin{enumerate}[resume]
\item Un desk qui juge une hausse peu probable achète 1\,000 contrats de décembre à
96.040. Quel est le prix final et son P\&L si le Comité maintient le taux~?
\item Et en cas de hausse~?
\item Quel est le P\&L espéré à la probabilité de 35~\% du marché~?
\item À quelle probabilité la transaction du desk atteint-elle le point mort, et pourquoi ce nombre est-il celui
du marché~?
\item Qu'est-ce qui, outre la décision, fait varier le règlement-livraison de décembre~?
\end{enumerate}

\medskip\noindent\textbf{Partie IV --- Jugement.}
\begin{enumerate}[resume]
\item Pourquoi «~une probabilité de 35~\% d'une hausse~» n'est-elle qu'une approximation, même sans passage~?
\item Pourquoi la probabilité risque-neutre pourrait-elle différer des projections
propres aux membres du Comité~?
\item Comment une hausse surprise en octobre modifierait-elle la lecture de décembre~?
\item Énoncer le \emph{résultat nommé}~: la probabilité implicite de hausse en décembre,
sans passage de fin d'année, puis avec un passage de 10 points de base.
\item En une phrase~: que représente le prix d'un contrat à terme sur taux à un mois~?
\end{enumerate}
\end{problem}

\section{Questions d'entretien}

\begin{interviewq}[$\star$ \iqroles{trader, chercheur}]\label{iq:m2:short-term-interest-rate-futures:1}
Comment lisez-vous la probabilité d'une hausse de taux à partir des prix des contrats à terme~?
\end{interviewq}

\begin{interviewq}[$\star$ \iqroles{trader, banque}]\label{iq:m2:short-term-interest-rate-futures:2}
Pourquoi les contrats à terme sur taux court cotent-ils $100$ moins le taux~?
\end{interviewq}

\begin{interviewq}[$\star\star$ \iqroles{chercheur, trader}]\label{iq:m2:short-term-interest-rate-futures:3}
Quel est l'\omterm{def:m2:short-term-interest-rate-futures:convexity}{ajustement de convexité} entre contrats à terme et taux à terme, et dans quel
sens agit-il~?
\end{interviewq}

\begin{interviewq}[$\star\star$ \iqroles{trader}]\label{iq:m2:short-term-interest-rate-futures:4}
Qu'est-ce qu'un pack, et pourquoi un trader négocierait-il le pack de la deuxième année plutôt que
les contrats les plus proches~?
\end{interviewq}

\begin{interviewq}[$\star\star$ \iqroles{développeur, chercheur}]\label{iq:m2:short-term-interest-rate-futures:5}
Vous construisez un service qui publie les taux implicites réunion par réunion
toutes les minutes. Qu'est-ce qui déraille en fin de mois, autour des réunions tenues tard dans le mois, et à
la fin d'année~?
\end{interviewq}

\begin{interviewq}[$\star\star\star$ \iqroles{chercheur}]\label{iq:m2:short-term-interest-rate-futures:6}
Les contrats à terme impliquent une probabilité de 35~\% d'une hausse. Est-ce la probabilité que vous devez utiliser
pour décider d'une transaction~? Que faudrait-il pour la convertir~?
\end{interviewq}
